\begin{keyinsight}[结论：基础安全边际有限；未来机会以风险调整期权进入合并价值]
德赛西威的投资逻辑建立在已确认的座舱、智能驾驶和网联收入，而不是单一芯片概念。2025年智能驾驶收入增长快于整体，但毛利率低于座舱且同比下行。当前价格已隐含约17倍公开2026年EPS；本院给出CNY86.4的基础业务多锚价值，再将2028年Physical AI的风险调整期权价值CNY2.49/股单列加入，合并参考值为CNY88.9、较CNY83.48约+6.5\%。期权不进入2026E EPS，且55\%权重仍对应截至FY2028未形成可衡量商业转化。行动为\stance{riskamber}{主业市场支持型观察 + 未来机会催化剂跟踪}。
\end{keyinsight}

\section{投资行动仪表盘}

本章同时回答两个不同问题：当前汽车电子盈利是否有足够安全边际，以及Physical AI能否形成新的2027--2028利润池。前者的答案是安全边际有限；后者已有产品、定点和技术外溢证据，但必须通过客户、交付、收入、毛利与现金的连续验证。评级标签描述研究行为，不是交易指令。

\begin{exhibitbox}[核心估值与行动]
\small
\begin{tabularx}{\textwidth}{L{2.4cm}Y Y Y Y}
\toprule
项目 & 结论 & 计算/证据 & 行动含义 & 失效条件 \\
\midrule
市场锚 & CNY83.48，2026-07-23盘中 & 总股本5.968亿股，市值CNY49.82bn，TTM P/E 21.36x & 只能作为盘中估值时点，非收盘价 & 价格、股本和估值倍数需在后续发布前刷新 \\
基础情景 & 收入CNY38.0bn，归母CNY2.80bn，EPS CNY4.69 & 2025A分部基础、2026Q1边界及公开盈利预测交叉检查 & 不是客户/车型拆分模型 & 若收入持续低增或智驾毛利率下行则下调 \\
基础业务价值 & CNY86.4，+3.5\% & 基本面CNY86.8权重80\%，市场锚CNY85.0权重20\%，Street 0\% & 不支持追高型核心配置 & 不把历史/无方法的券商目标价加权 \\
Physical AI期权 & CNY2.49/股 & FY2028三档结果以12\%折现；55\%/35\%/10\%主观风险权重 & 进入合并价值，不进入基础EPS & 若2026后仍无交付、收入或经济学证据，则下调而非延长期权 \\
合并参考值与行动 & CNY88.9，+6.5\%；市场支持型观察 + 催化剂跟踪 & CNY86.4基础价值 + CNY2.49风险调整期权 & 等待中报/量产/收入证据或更好的安全边际 & 主业业绩、现金和客户链未达阈值则转观察名单 \\
\bottomrule
\end{tabularx}
\sourcenote{公司2025年年报、2026Q1报告、2026-07-23市场快照；本院情景估值。}
\end{exhibitbox}

该估值并不把“合并参考值略高于现价”误读为确定买入信号，也不把它误读为未来机会已经结束。市场锚之所以保留20\%权重，是因为成长型 Tier-1 的交易价格会反映已经可见的客户覆盖、平台进展和行业情绪；但当前没有可审计的券商目标价可作为第三个价格锚，且客户/车型/ASP没有披露到可做精确增量EPS的程度。因此，基础估值的中心是交付质量；未来机会的中心是下一个可以改变概率的商业化事实。

\section{回报归因与上、下行触发}

基础情景的上行主要来自收入重新回到中双位数增长和智能驾驶毛利率止跌，而非估值倍数的大幅扩张。牛市情景需要高阶座舱、智驾量产和新增平台收入三者同时证明；熊市情景则反映交付放缓、成本/费用压力或营运资本恶化。它们是一组概率条件，不是对单季度利润的年化。Physical AI的远期期权不藏入这三档：先以客户--交付--收入--毛利--现金五步升级，再建立单列的2027--2028情景。

\begin{exhibitbox}[三档价值与验证条件]
\small
\begin{tabularx}{\textwidth}{L{1.8cm}Y Y Y Y}
\toprule
情景 & 2026E收入/归母/EPS & 估值方法 & 价值 & 必须发生的事情 \\
\midrule
Bear & CNY35.5bn / CNY2.50bn / CNY4.19 & 16.5x P/E & CNY69.1 & 收入放缓、智驾毛利率继续承压、营运资本占用上升 \\
Base & CNY38.0bn / CNY2.80bn / CNY4.69 & 18.5x P/E & CNY86.8 & 收入恢复中双位数、智驾毛利不继续恶化、现金转换正常 \\
Bull & CNY40.0bn / CNY3.15bn / CNY5.28 & 20.0x P/E & CNY105.6 & 量产及产品结构上移得到客户、收入和利润的可审计验证 \\
\bottomrule
\end{tabularx}
\sourcenote{本院情景；价格为CNY83.48盘中锚。}
\end{exhibitbox}

对投资委员会而言，最重要的纪律是把“订单年化销售额”和“已确认收入”分开。公司披露2025年新项目订单年化销售额超过CNY35.0bn，其中座舱超过CNY20bn、智能驾驶超过CNY13bn；其定义是订单全生命周期销售额除以生命周期年数。它提供需求可见度，却不能跳过开发、SOP、交付、客户验收和收入确认。基准EPS没有把这一指标直接转化为收入。

\section{排名、风险预算与监测行为}

本报告覆盖池只对德赛西威估值：华阳集团、均胜电子、科博达、经纬恒润和全球Tier-1用于竞争参照，OEM与芯片平台是需求/供应链锚点，均不被当作本公司收入证明。行动排序使用四个权重：已确认经营事实40\%、客户链与量产证据25\%、估值安全边际20\%、近季验证可得性15\%。在这个框架下，德赛西威具备前两项优势，但安全边际和近季盈利验证均为中等。对Physical AI另加一条研究线：能否在主业现金纪律不恶化的前提下，先后获得具名商业身份、实际交付和收入/毛利证据。

\begin{riskbox}[监测纪律]
在正式中报前，不将本报告等同于建仓建议。基础升级需要收入增速恢复、智能驾驶毛利率企稳、经营现金流与应收/存货相协调；Physical AI升级至少需要具名客户/场景、实际交付或收入中的一项，再观察毛利和现金。降级触发包括收入持续低增、智能驾驶毛利率再降、营运资本明显快于销售，或新增平台仍只有订单/定点而无收入证据。
\end{riskbox}
